\chapter{Задача: гидриды в стенке трубы}\label{ch:problem}

\section{Откуда в трубе гидриды}

Оболочки твэлов и направляющие каналы тепловыделяющих сборок делают из сплавов циркония:
Zircaloy-4, Zr--1Nb, Э110, Э635 (Zr--Nb--Sn--Fe). В реакторе цирконий окисляется водой, и часть
водорода, освободившегося при коррозии, уходит в металл. При рабочей температуре водород
растворён, но растворимость быстро падает с температурой: при 400\,$^\circ$C металл держит
около 170 частей на миллион по массе (ppm), при комнатной --- меньше одной. Всё лишнее при
остывании выпадает в виде \emph{гидрида} --- отдельной фазы $\delta$-ZrH$_{1.66}$, у которой
на атом циркония приходится 1.66 атома водорода.

Гидрид выпадает тонкими пластинками: толщина --- доли микрона, длина --- единицы--десятки
микрон. Он хрупкий, и трещина охотно идёт вдоль пластинок. Поэтому важно, \emph{как они лежат}.

\fig{F01_Tube}{Слева --- сечение трубы. В нём пластинки гидрида обычно лежат вдоль окружности
(синие). Если труба остывает под окружным напряжением $\sigma_\theta$ (внутреннее давление),
новые пластинки могут лечь по радиусу (красные). Такие гидриды выстраиваются поперёк стенки и
опасны: трещине почти не нужно идти по металлу.}

\begin{idea}
Окружные гидриды лежат «вдоль слоёв» стенки, и трещина, идущая изнутри наружу, вынуждена
пересекать их по одному. Радиальные гидриды лежат поперёк стенки --- по пути трещины. Поэтому
переориентация гидридов под нагрузкой --- главный вопрос при сухом хранении отработавшего
топлива и при испытаниях колец, которыми занимается проект.
\end{idea}

\section{Кристалл, текстура и почему гидриды обычно окружные}

Цирконий --- гексагональный металл. У каждого зерна есть ось $c$ и перпендикулярная ей
\emph{базисная} плоскость. Гидрид выпадает пластинкой почти параллельно базисной плоскости
(точнее, на плоскостях типа $\{10\bar17\}$, в $14.7^\circ$ от базиса). Значит, ориентацию
пластинки задаёт ориентация зерна.

При прокатке трубы зёрна поворачиваются так, что оси $c$ смотрят в основном по радиусу с наклоном
$\pm\chi_0\approx 30$--$40^\circ$ к окружности. Это \emph{текстура}, и её обычно описывают
факторами Кернса $f_r$, $f_\theta$, $f_z$ --- средними квадратами косинусов оси $c$ с тремя
направлениями ($f_r+f_\theta+f_z=1$). У листа Zircaloy-4 из опытов Lepine $f_r\approx0.59$,
$f_\theta\approx0.30$, $f_z\approx0.12$.

Если ось $c$ смотрит почти по радиусу, базисная плоскость --- а с ней и пластинка --- лежит почти
по окружности. Отсюда окружные гидриды «по умолчанию». Чтобы получить радиальный гидрид, нужно
либо зерно с осью $c$ по окружности (таких меньшинство), либо механизм, который выстраивает
пластинки в радиальную цепочку. Оба пути есть в модели; второй, коллективный, оказался
важнее, чем казалось (гл.~\ref{ch:results}).

\begin{exercise}
Проверьте связь факторов Кернса с наклоном. Пусть оси $c$ лежат в плоскости $r$--$\theta$ под
углом $\chi$ к радиусу, равновероятно $\pm\chi$. Покажите, что $f_\theta/(f_r+f_\theta)=\sin^2\chi$.
Найдите $\chi$ для листа Lepine ($f_r=0.59$, $f_\theta=0.30$).
\end{exercise}

\section{Сколько гидрида и сколько пластинок}

Водорода немного, а гидрида по объёму --- ещё меньше процента--двух. Пусть в металле $H$~ppm
водорода и весь он выпал. В гидриде ZrH$_{1.66}$ на 1 кг циркония приходится
$1.66\cdot1.008/91.22\cdot10^6\approx18\,000$~ppm водорода, а гидрид на $\approx15\%$ объёмнее
металла с тем же цирконием. Значит, доля объёма гидрида
\begin{equation}
f_{\text{гидр}}\approx\frac{H}{18\,013}\cdot1.15=\frac{H}{15\,660}.
\label{eq:frac}
\end{equation}
Число $C_{\text{гидр}}=15\,660$~ppm будет встречаться дальше: это «сколько водорода сидит в
единице площади гидрида», если мерить в ppm металла.

\begin{exercise}
Какая доля сечения занята гидридом при 178 ppm (опыты Lepine)? Сколько пластинок
$5\times0.6$~мкм поместится в поле $240\times240$~мкм, если весь гидрид выпал такими пластинками?
\end{exercise}

\section{Что нужно объяснить}

Переориентацию изучают так: образец насыщают водородом, нагревают до $T_{\max}$ (часть или весь
гидрид растворяется), нагружают и медленно охлаждают под нагрузкой. Потом на шлифе считают
долю радиальных гидридов (RHF, гл.~\ref{ch:measures}). Опыты дают несколько закономерностей,
которые хорошая модель должна воспроизводить, получая на вход только условия опыта:

\begin{enumerate}
\item \textbf{Порог.} Ниже некоторого напряжения радиальных гидридов почти нет, выше --- их доля
  быстро растёт. У холоднодеформированного Zircaloy-4 порог $\approx150$--$175$~МПа
  (Cinbiz 2016, Desquines 2014), у рекристаллизованных сплавов и Zr--Nb --- 75--100~МПа
  (Sakamoto 2006, Colas 2010, Son 2026).
\item \textbf{Порог почти не зависит от водорода} при 100--450 ppm (полка Desquines), но при малом
  водороде радиальных больше даже без нагрузки (Son 2026).
\item \textbf{Форма перехода:} у Cinbiz --- ступенька за 15--30~МПа, у Son --- плавный рост на 60~МПа.
\item \textbf{Неполное растворение:} если $T_{\max}$ ниже температуры полного растворения,
  оставшиеся окружные гидриды служат зародышами, и радиальных меньше.
\item \textbf{Межзёренные гидриды:} в мелкозернистом холоднодеформированном металле 92\% длины
  гидридов лежит на границах зёрен (Son 2026).
\item \textbf{Скорость охлаждения:} медленнее --- больше и длиннее радиальных (Min 2013).
\item \textbf{Двухосность:} при одновременном осевом растяжении порог ниже (Cinbiz 2016).
\item \textbf{Морфология:} крупный радиальный гидрид при большом увеличении оказывается
  ступенчатой стопкой мелких пластинок (Puls 2012, гл.~3).
\end{enumerate}

Не все эти пункты --- честная проверка для двумерной модели (гл.~\ref{ch:limits}), и не все
модель сейчас воспроизводит. Это и есть содержание глав~\ref{ch:results}--\ref{ch:limits}.

\section{Подход модели в одном абзаце}

Берём поперечное сечение стенки ($r$--$\theta$) и покрываем его сеткой клеток 0.4~мкм. Зёрна ---
многоугольники Вороного со своей ориентацией. Температура падает с заданной скоростью. На каждом
шаге по времени в каждой клетке считается, насколько выгодно здесь зародиться пластинке:
это зависит от пересыщения раствора (гл.~\ref{ch:solubility}), от напряжений, которые создают
уже выросшие пластинки и нагрузка (гл.~\ref{ch:misfit}--\ref{ch:plasticity}), и от теории
зарождения (гл.~\ref{ch:nucleation}). Зародыши разыгрываются случайно, растут, забирают водород,
водород диффундирует (гл.~\ref{ch:hydrogen}). Так шаг за шагом (гл.~\ref{ch:automaton}) из
физики получается структура, которую можно мерить так же, как шлиф (гл.~\ref{ch:measures}).

\begin{idea}[Философия]
В модели должно быть как можно меньше подогнанных чисел. Каждый параметр либо измерен, либо
рассчитан, либо оценён из теории; подогнанные помечены, и их число --- мера того, насколько
модель ещё «кривая». Сравнение с опытом --- не «попасть в три статьи одним набором чисел», а
получить на вход условия каждого опыта и выдать его результат, калибруясь на одном наборе данных.
\end{idea}
